\begin{document}
\maketitle

\begin{frame}
  \titlepage
\end{frame}

\cleardoublepage

\tableofcontents

\cleardoublepage

\section{Einführung}

Systemd ist ein umfangreiches Management-Tool für System und Services mit
entsprechenden Vor- und Nachteilen. SysV Init ist nur für ein kleines Subset verantwortlich.

\begin{frame}{Systemd}
  \begin{itemize}
    \item All-In-One
    \item Abhängigkeiten
    \item Parallele Ausführung
    \item Geschwindigkeit
    \item Prozesse überwachen
    \item Ressourcen + Security
  \end{itemize}
\end{frame}

\section{Bestandteile}

\begin{frame}{Bestandteile}
  \begin{itemize}
\item Services
\item Logging
\item Hostname/Locale/Keyboard/Time/...
\item Temporäre Dateien
\item Timers (Cron-Ersatz)
\item ...
\end{itemize}
\end{frame}

\subsection{Unit-Typen}

\begin{frame}{Units}
  \begin{itemize}
  \item Services (.service)
  \item Gruppen (.target)
  \item Slice (.slice)
  \item Scope (.scope)
  \item Dateisysteme (.mount / .automount)
  \item Überwachung (.path)
  \item Sockets (.socket)
  \item Netzwerke (.network)
  \item Device (.device)
  \item ...
\end{itemize}
\end{frame}

Der Unit-Typ Slice erzeugt die Cgroup-Hierarchie für Resource-Management.

Der Unit-Typ Scope gruppiert die Worker-Prozesse für einen Service.

\subsection{Targets}
Targets sind einfach eine Gruppe von Unitfiles.

Außerdem gibt es Targets, die beim Bootvorgang aufgerufen werden und grob
die Funktion von Runlevels ausüben:

\begin{frame}[fragile]{Targets}
  \begin{itemize}
  \item graphical.target
  \item multiuser.target
  \item rescue.target
  \item default.target
  \end{itemize}
\end{frame}

Dabei entspricht graphical grob dem Runlevel 1, multiuser dem Runlevel 3,
rescue dem Runlevel 1 und default ist ein Alias für das aktive Target
beim Booten.

\section{Kommandozeile}

\begin{frame}{Kommandozeile}
  \begin{itemize}
  \item systemctl
  \item journalctl
  \end{itemize}
\end{frame}

\subsection{Services}

\begin{frame}[fragile]{Services}
 \begin{lstlisting}
    systemctl status nginx
    systemctl stop nginx
    systemctl start nginx
    systemctl enable nginx
    systemctl disable nginx
  \end{lstlisting}
\end{frame}

\begin{frame}[fragile]{Service Status}
\begin{lstlisting}
~# systemctl status ntp
ntp.service - LSB: Start NTP daemon
   Loaded: loaded (/etc/init.d/ntp; generated; vendor preset: enabled)
   Active: active (running) since Mon 2018-03-26 10:15:10 CEST; 25min ago
     Docs: man:systemd-sysv-generator(8)
  Process: 581 ExecStart=/etc/init.d/ntp start (code=exited, status=0/SUCCESS)
    Tasks: 2 (limit: 4915)
   CGroup: /system.slice/ntp.service
           594 /usr/sbin/ntpd -p /var/run/ntpd.pid -g -u 108:113
\end{lstlisting}
\end{frame}

\subsection{Fehlerbehandlung}

systemctl start gibt keine Fehlermeldungen aus. Im Default-Modus kann es sogar dazu kommen, daß gar keine Meldung kommt.

\begin{frame}[fragile]{Fehlerbehandlung}
\begin{lstlisting}
~# systemctl start nginx
Job for nginx.service failed because the control process
exited with error code.
See "systemctl status nginx.service" and "journalctl -xe" 
for details.
\end{lstlisting}
\end{frame}

\begin{frame}[fragile]{Fehlerbehandlung}
\begin{lstlisting}
~# systemctl start elasticsearch
~#
\end{lstlisting}
\end{frame}

\begin{frame}[fragile]{Fehlerbehandlung}
\begin{lstlisting}
~# systemctl --failed
UNIT                      LOAD   ACTIVE SUB    DESCRIPTION
elasticsearch.service loaded failed failed Elasticsearch
\end{lstlisting}
\end{frame}

\begin{frame}{Fehlerbehandlung}
  \begin{figure}[!ht]
     \centering
     \includegraphics[width=0.9\linewidth]{img/stress.jpg}
  \end{figure}
\end{frame}

\begin{frame}[fragile]{Fehlerbehandlung}
\begin{lstlisting}
# There is insufficient memory for the Java Runtime Environment to continue.
# Native memory allocation (mmap) failed to map 3046768640 bytes for committing reserved memory.
\end{lstlisting}
\end{frame}

\subsection{Units}

Mit SysV Init finden sich die Skripte für die Services in \verb|/etc/init.d|.

\begin{frame}[fragile]{Unit anzeigen}
\begin{lstlisting}
~# systemctl cat nginx
# /lib/systemd/system/nginx.service
[Unit]
After=network.target

[Service]
Type=forking
PIDFile=/run/nginx.pid
ExecStartPre=/usr/sbin/nginx -t -q -g 'daemon on; master_process on;'
ExecStart=/usr/sbin/nginx -g 'daemon on; master_process on;'
ExecReload=/usr/sbin/nginx -g 'daemon on; master_process on;' -s reload
ExecStop=-/sbin/start-stop-daemon --quiet --stop --retry QUIT/5 --pidfile /run/nginx.pid
TimeoutStopSec=5
KillMode=mixed

[Install]
WantedBy=multi-user.target
\end{lstlisting}
\end{frame}

\begin{frame}[fragile]{Liste der Units}
  \begin{itemize}
    \item
Units auflisten:

    \verb|systemctl list-units|
    \item
Unitdateien:

\verb|systemctl list-unit-files|
  \end{itemize}
\end{frame}

\subsection{Filter}

\begin{frame}[fragile]{Filter}
  \begin{itemize}
    \item
	    Unit-Typ (Services):

    \verb|systemctl list-units --type=service|

\item
	Inaktive und fehlende Services:

    \verb|systemctl list-units --type=service --all|
  \end{itemize}
\end{frame}

\subsection{Status Unitdateien}

\begin{frame}[fragile]{Status Unitdateien}
\begin{itemize}
\item enabled
\item disabled
\item generated
\item masked
\item static
\end{itemize}
\end{frame}


\subsection{Targets}

\begin{frame}[fragile]{Abhängigkeiten anzeigen}
\begin{lstlisting}
systemctl list-dependencies multi-user.target
\end{lstlisting}
\end{frame}

\begin{frame}[fragile]{Default Target}
  \begin{itemize}
    \item Anzeige

    \verb|systemctl get-default|

    \item Voreinstellung

    \verb|systemctl set-default multi-user.target|

    \item Laufzeitänderung

    \verb|systemctl isolate multi-user.target|
  \end{itemize}
\end{frame}

\section{Unitdateien}

\subsection{Dateisystem}

\begin{frame}[fragile]{Dateisystem}
  \begin{itemize}
\item Standard

	\verb|/lib/systemd/system|
\item Custom (Lokal, Ansible)

	\verb|/etc/systemd/system|

\item Runtime

	\verb|/run/systemd/system|
\end{itemize}
\end{frame}

\subsubsection{target.wants}

\subsection{Struktur Unitdatei}

\begin{frame}[fragile]{Struktur Unitdatei}
 \begin{lstlisting}
[Unit]
Description=The PHP FastCGI Process Manager
After=network.target

[Service]
Type=notify
PIDFile=/var/run/php5-fpm.pid
ExecStartPre=/usr/lib/php5/php5-fpm-checkconf
ExecStart=/usr/sbin/php5-fpm --nodaemonize \
    --fpm-config /etc/php5/fpm/php-fpm.conf
ExecReload=/bin/kill -USR2 $MAINPID

[Install]
WantedBy=multi-user.target
\end{lstlisting}
\end{frame}

\subsection{Direktiven}

\begin{frame}{Direktiven}
% https://pixabay.com/de/fische-meer-ozean-fischschwarm-22646/
  \begin{figure}[!ht]
     \centering
     \includegraphics[width=0.8\linewidth]{img/fish-22646_1920.jpg}
  \end{figure}
\end{frame}

\begin{frame}[fragile]{Abschnitte}
	\begin{itemize}
		\item \verb|[Unit]|
		\item \verb|[Service]|
		\item \verb|[Install]|
	        \item ...

	\end{itemize}	
\end{frame}

\begin{frame}[fragile]{Environment}
\begin{lstlisting}
[Service]	
Environment=ES_HOME=/usr/share/elasticsearch
Environment=CONF_DIR=/etc/elasticsearch/es1

EnvironmentFile=-/etc/default/es1_elasticsearch

WorkingDirectory=/usr/share/elasticsearch

User=elasticsearch
Group=elasticsearch
\end{lstlisting}
\end{frame}

\begin{frame}[fragile]{Restart}
\begin{lstlisting}
[Service]	
Restart=always
\end{lstlisting}
\end{frame}

\subsection{Units bearbeiten}
\begin{frame}{Units bearbeiten}
  \begin{itemize}
    \item Unitdatei kopieren und bearbeiten
    \item Drop-in erstellen
  \end{itemize}
\end{frame}

\begin{frame}[fragile]{Drop-in erstellen}
	\begin{lstlisting}
# systemctl show -p Restart nginx
Restart=no

# systemctl edit nginx
[Service]
Restart=always

# systemctl show -p Restart nginx
Restart=always
\end{lstlisting}
\end{frame}

Hierbei gilt es auf die korrekte Syntax zu beachten, z.B. \verb|[service]| statt \verb|[Service]|
führt zur Mißachtung der Anweisungen.

\subsection{Service-Typen}

\begin{frame}{Service-Typen}
\begin{itemize}
\item simple
\item notify
\item forking
\item ...
\end{itemize}
\end{frame}

\begin{frame}[fragile]{Typ simple}
\begin{lstlisting}
ExecStart=/usr/share/elasticsearch/bin/elasticsearch \
                      -p ${PID_DIR}/elasticsearch.pid \
                      --quiet
\end{lstlisting}
\end{frame}

\begin{frame}{Fire and Forget}
% https://pixabay.com/de/feuerwerk-pyrotechnik-benutzt-1156808/	
  \begin{figure}[!ht]
      \centering
      \includegraphics[width=0.9\linewidth]{img/fireworks-1156808_1920.jpg}
   \end{figure}
\end{frame}

\begin{frame}[fragile]{Workarounds}
  \begin{itemize}
     \item ExecStartPre

\verb|ExecStartPre=/usr/lib/php5/php5-fpm-checkconf|

     \item Typ notify

	Änderung am Service erforderlich
  \end{itemize}
\end{frame}

\section{Protokollierung}

journalctl ist für die Protokollierung zuständig. Die Daten werden in einem Binärformat gespeichert.

\begin{frame}[fragile]{Protokollierung}
\begin{itemize}
\item Binärformat
\begin{itemize}
\item \verb|/run/log/journal| (gelöscht beim Booten)
\item \verb|/var/log/journal| (persistent)
\end{itemize}
\item Syslog
	\begin{itemize}
	\item rsyslog liest Journal
	\item \verb|ForwardToSyslog=yes|
	\end{itemize}
\end{itemize}
\end{frame}

\begin{frame}[fragile]{Protokollierung}
  \begin{lstlisting}
    journalctl -f
    \end{lstlisting}
\end{frame}
  
\subsection{Filter}

\begin{frame}[fragile]{Filter für journalctl}
  \begin{itemize}
  \item Service

	  \verb|journalctl -u nginx|
  \item Zeit

	  \verb|journalctl --since 09:00 --until "1 hour ago"|
  \item Priorität

	  \verb|journalctl -p err|
  \end{itemize}
\end{frame}

\section{Systemd \& Sysvinit}

\begin{frame}{Systemd \& SysV-Init}
  \begin{itemize}
  \item Initskripte
  \item Runlevel
  \item Inetd
  \end{itemize}
\end{frame}

\subsection{Einbindung alte Services}

\begin{frame}[fragile]{Einbindung alte Services}
  \begin{itemize}
    \item Systemstart
    \item Generierung Units

	    \verb|systemd-sysv-generator|
  \end{itemize}
\end{frame}

\begin{frame}[fragile]{NTP Service I}
\begin{lstlisting}
~# systemctl cat ntp
# /run/systemd/generator.late/ntp.service
# Automatically generated by systemd-sysv-generator

[Unit]
Documentation=man:systemd-sysv-generator(8)
SourcePath=/etc/init.d/ntp
Description=LSB: Start NTP daemon
Before=multi-user.target
Before=multi-user.target
Before=multi-user.target
Before=graphical.target
After=network-online.target
After=remote-fs.target
Wants=network-online.target
\end{lstlisting}
\end{frame}

\begin{frame}[fragile]{NTP Service II}
\begin{lstlisting}
[Service]
Type=forking
Restart=no
TimeoutSec=5min
IgnoreSIGPIPE=no
KillMode=process
GuessMainPID=no
RemainAfterExit=yes
SuccessExitStatus=5 6
ExecStart=/etc/init.d/ntp start
ExecStop=/etc/init.d/ntp stop
\end{lstlisting}
\end{frame}

\subsection{Runlevel}

\begin{frame}{Runlevel}
\begin{itemize}
\item Obsolet in Systemd
\item Vergleichbar mit Targets
\begin{description}[0]
\item[1] emergency.target
\item[3] multi-user.target
\item[5] graphical.target	
\end{description}		
\end{itemize}	
\end{frame}

\section{Abschluß}
\subsection{Vor- und Nachteile}

\begin{frame}{Vor- und Nachteile}
\begin{itemize}
\item[\goodsmile] Abhängigkeiten
\item[\goodsmile] Geschwindigkeit
\item[\neutralsmile] Unix-Philosophie	
\item[\sadsmile] Fallgrube Service-Typ "simple"
\item[\sadsmile] Dokumentation
\end{itemize}
\end{frame}

\subsection{Referenzen}

\begin{frame}{Referenzen}
\begin{itemize}	
	\item \href{https://wiki.archlinux.org/index.php/systemd}{ArchWiki}

	\item \href{https://www.redhat.com/files/summit/session-assets/2017/S103870-Demystifying-systemd.pdf}{Demystifying systemd}
	\item \href{https://www.digitalocean.com/community/tutorials/understanding-systemd-units-and-unit-files}{https://www.digitalocean.com/community/tutorials/understanding-systemd-units-and-unit-files}

	\item \href{https://www.digitalocean.com/community/tutorials/how-to-use-journalctl-to-view-and-manipulate-systemd-logs}{https://www.digitalocean.com/community/tutorials/how-to-use-journalctl-to-view-and-manipulate-systemd-logs}
\end{itemize}
\end{frame}

\end{document}

%%% Local Variables: 
%%% mode: latex
%%% TeX-master: t
%%% End: 
